\SourceLeaf{282}%
ein Minimum werden soll, wofür nach den Regeln der Variationsrechnung sich die Gleichung ergibt
\[
\frac{\mathrm{d}n}{\mathrm{d}u}\mathord{.}\BareRoot{(1+\frac{\mathrm{d}u^2}{\mathrm{d}x^2})}\mathrm{d}x=\mathrm{d}\frac{\frac{n\mathrm{d}u}{\mathrm{d}x}}{\BareRoot{(1+\frac{\mathrm{d}u^2}{\mathrm{d}x^2})}}
\]
oder
\[
\frac{\mathrm{d}n}{\mathrm{d}u}\mathord{.}\frac{\mathrm{d}x}{\cos\psi}=\mathrm{d}\mathord{.}n\sin\psi
\]
Unter \(\tfrac{\mathrm{d}n}{\mathrm{d}u}\) ist der partielle Differentialquotient verstanden. Diese Formel ist strenge und allgemeingültig. Für unsern Zweck aber, wo bloss das zwischen \(F\) und \(G\) liegende Stück der Curve \(M\) in Betracht kommt, in deren sämmtlichen Punkten \(u\) und \(\psi\) nur sehr kleine Werthe haben können, dürfen wir \(1\) anstatt \(\cos\psi\) und \(\operatorname{tang}\psi\) anstatt \(\sin\psi\) schreiben, mithin
\[
\frac{\mathrm{d}n}{\mathrm{d}u}\mathord{.}\mathrm{d}x=\mathrm{d}\mathord{.}n\operatorname{tang}\psi
\]
oder
\[
n\operatorname{tang}\psi=\int\frac{\mathrm{d}n}{\mathrm{d}u}\mathrm{d}x+\text{Const.}
\]
setzen, zugleich aber auch in dieser Formel anstatt der Werthe, welche \(n\) und \(\tfrac{\mathrm{d}n}{\mathrm{d}u}\) in der Linie \(M\) haben, diejenigen anwenden, welche in den correspondirenden Punkten der Linie \(N\) (für \(u=0\) oder \(y=0\)) Statt finden, und folglich mit den Werthen von \(\tfrac{1}{m}\) und \(-\tfrac{\mathrm{d}m}{mm\mathrm{d}u}=-\tfrac{\mathrm{d}m}{mm\mathrm{d}y}\) übereinstimmen.

Zur bequemern Ausführung der weitern Entwicklungen sollen jetzt die Abscissen von dem Punkte \(F\) an gezählt, oder in diesem Punkte \(x=0\), in \(G\) hingegen \(x=h\) gesetzt werden; ich setze ferner \(\tfrac{\mathrm{d}m}{m\mathrm{d}y}=l\), welches im Allgemeinen zwar Function von \(x\) und \(y\) ist, hier aber bloss nach seinem in der Linie \(N\) oder für \(y=0\) geltenden Werthe, also als Function von \(x\) allein betrachtet wird; endlich seien \(\psi^0,m^0,l^0\), die bestimmten Werthe von \(\psi,m,l\) in dem Punkte \(F\), und \(\psi',m',l'\) die in dem Punkte \(G\). Die obige Formel wird hienach
\[
\operatorname{tang}\psi=\frac{m\operatorname{tang}\psi^0}{m^0}-m\int\frac{l}{m}\mathrm{d}x
\]
wo die Integration von \(x=0\) anfängt. Nehmen wir nun an, dass \(l\) und \(m\) in folgende nach Potenzen von \(x\) fortschreitende Reihen
